\chapter{Functions}
\label{chap:functions}

The programs of the previous chapter fit in \token{main}, and it was already
getting crowded. A program that draws two triangles writes the loops that draw a
triangle twice. A program that tests whether several numbers are prime writes
the primality test once for each. Each copy is one more place to read, and one
more place to forget when the code has to change.

A cookbook does not repeat the recipe of the pastry in every pie that uses it. It
writes the recipe once, gives it a name, and every pie says ``make the pastry of
page 12''. A \textit{function} is the same thing in a program: a piece of code
with a name, written once, that the rest of the program runs by calling that
name. Chapter~\ref{chap:basics} already declared a few, without parameters or
with a line of explanation. This chapter shows how to give a function inputs and
a result, how a function can call itself, and how several functions can share a
name. It ends with a program large enough to need all of that: a calendar.

\notebox{From this chapter on, the complete programs of the course are also
  distributed as files, in the archive \textit{ymir-book-examples} published
  with each release of the book, at
  \url{https://github.com/GNU-Ymir/specifications/releases}. The caption of a
  listing taken from one of them gives its path in the archive, such as
  \textit{examples/functions/calendar.yr}. Each file is a whole program, compiled
  with \token{gyc} like any other
  (Section~\ref{sec:basics:compilation_execution}).}

Part II covers functions exhaustively, in
Section~\ref{sec:global:functions}.

\minitoc%

\vfill\pagebreak
\section{Declaring and calling a function}
\label{sec:functions:declaring}

The report below prints a line of dashes three times. Rather than writing the
same \token{println} three times, it declares a function \token{separator} that
prints the line, and calls it wherever a line is needed.

\lstinputlisting[style=coloredverbatim, caption={A function called three times, \textit{examples/functions/separator.yr}}, label=lst:functions:separator]{examples/functions/separator.yr}
\vspace{-10pt}%
\begin{lstlisting}[style=bashVerb, escapechar=@+]
@+\textcolor{teal!80}{alice@dev:\mtilde}@+$ gyc separator.yr -o separator
@+\textcolor{teal!80}{alice@dev:\mtilde}@+$ ./separator
--------------------
  Monthly report
--------------------
Sales: 1250
Costs:  830
--------------------
\end{lstlisting}

A function declaration has four parts, which lines~3 to~5 show in order:

\begin{itemize}
\item the keyword \token{fn}, which announces a function;
\item its name, \token{separator}, chosen by the programmer under the rules of
  identifiers (Section~\ref{sec:types:constants_variables}). By convention the name
  starts with a lowercase letter, and a name of several words capitalizes each
  word after the first, as in \token{printCalendar};
\item a pair of parentheses, which holds the function's parameters -- here there
  are none, the next subsection gives it some;
\item its \textit{body}, a block (Section~\ref{sec:control:blocks}) that holds
  the instructions to run.
\end{itemize}

A declaration runs nothing by itself. The body runs when the function is
\textit{called}, by writing its name followed by parentheses, as on lines~8,
10 and 13. A call jumps to the start of the body, runs it to the end, and then
comes back to the caller, just after the call, where the program carries on.
Figure~\ref{fig:functions:call_return} follows the cursor through the first two
calls. \token{main} itself is a function, the one the program calls first, and
\token{println} is a function too, declared in the module \token{std::io}.

\begin{figure}[h]
  \centering
  \begin{adjustbox}{max width=\linewidth}
    \begin{tikzpicture}[node distance=6mm]

      \node[cfterm] (start) {start of main};
      \node[cfstep, below=of start] (call1) {separator();};
      \node[cfstep, below=12mm of call1] (title) {println("  Monthly report");};
      \node[cfstep, below=of title] (call2) {separator();};
      \node[cfstep, below=12mm of call2] (sales) {println("Sales: 1250");};
      \node[cflab, below=4mm of sales] (more) {\dots};

      \node[cfstep, right=22mm of call1, yshift=-9mm] (body1) {println("-----...");};
      \node[cfstep, right=22mm of call2, yshift=-9mm] (body2) {println("-----...");};
      \node[cflab, right=2mm of body1, align=left] {body of\\separator};
      \node[cflab, right=2mm of body2, align=left] {body of\\separator};

      \draw[cfflow] (start) -- (call1);
      \draw[cfflow] (call1.east) -| node[cflab, above, pos=0.25] {call} (body1.north);
      \draw[cfflow] (body1.south) |- node[cflab, below, pos=0.75] {return} (title.east);
      \draw[cfflow] (title) -- (call2);
      \draw[cfflow] (call2.east) -| node[cflab, above, pos=0.25] {call} (body2.north);
      \draw[cfflow] (body2.south) |- node[cflab, below, pos=0.75] {return} (sales.east);
      \draw[cfflow] (sales) -- (more);
    \end{tikzpicture}
  \end{adjustbox}
  \caption{\label{fig:functions:call_return} The first two calls of
    Listing~\ref{lst:functions:separator}. Each call runs the body of
    \token{separator} from the start, then control returns to the statement that
    follows the call.}
\end{figure}


If the separator should become a line of stars, one line of the program changes,
line~4, and all three separators follow. With three copies of the
\token{println}, all three would have to be found and changed, and missing one
would leave a report whose lines do not match.

As Chapter~\ref{chap:basics} said, the order in which functions are declared in
the file does not matter: \token{main} may call a function declared below it.

\subsection{Parameters}
\label{sec:functions:parameters}

A function that always does exactly the same thing is of limited use. A function
drawing a line of stars is more useful if each call can say how long the line
is. A \textit{parameter} is a variable of the function whose value is given by
each call. It is declared between the parentheses of the declaration, as a name
followed by a colon and a type, like a variable declaration without its
\token{let}. The value that a call gives it, written between the parentheses of
the call, is called an \textit{argument}.

\lstinputlisting[style=coloredverbatim, caption={Functions with parameters, \textit{examples/functions/rectangle.yr}}, label=lst:functions:rectangle]{examples/functions/rectangle.yr}
\vspace{-10pt}%
\begin{lstlisting}[style=bashVerb]
************

******
******
******
\end{lstlisting}

\token{line} has one parameter, \token{width}, of type \token{i32}. The call
\token{line(12)} on line~17 runs the body of \token{line} with \token{width}
holding 12, and each call gets the value it passes: \token{width} is 6 during the
calls that \token{rectangle} makes.

\token{rectangle} has two parameters, separated by a comma, and a call gives
them two arguments in the same order: in \token{rectangle(6, 3)}, \token{width}
is 6 and \token{height} is 3. It also shows that a function can call another
function. \token{rectangle} does not know how a line is drawn, only that
\token{line} draws one, so a rectangle is three lines of 6 stars.

An argument can be any expression of the right type, not only a literal:
\token{line(n)}, \token{line(n * 2 + 1)} or \token{line(read!i32())} are all
calls to \token{line}. The expression is computed first, then its value is given
to the parameter.

\subsection{Parameters are constants}

A parameter holds a copy of its argument, and the function cannot change it:
like a variable declared without \token{mut}
(Section~\ref{sec:types:variable_declaration}), a parameter is a constant. The
next listing tries to count down with its parameter, and is refused.

\begin{lstlisting}[style=coloredverbatimError, escapechar=@, caption=Assigning a parameter]
use std::io;

fn countdown(from: i32) {
  while from > 0 {
    println(from);
    @\hb{from}@ -= 1; // error
  }
}

fn main() {
  countdown(3);
}
\end{lstlisting}

The compiler says \errmsg{E4156}{from is a value parameter, and cannot be used as
  a lvalue},
an \textit{lvalue} being something that can stand on the left of an assignment.
A function that needs a value it can change declares a variable of its own,
initialized from the parameter.

\begin{lstlisting}[style=coloredverbatim, escapechar=@, caption=A mutable copy of a parameter]
use std::io;

fn countdown(from: i32) {
  @\hcb{let mut n = from;}@
  while n > 0 {
    println(n);
    n -= 1;
  }
}

fn main() {
  countdown(3);
}
\end{lstlisting}

A parameter follows the other rules of variables, too. One that the body never
reads is refused, as an unused variable is
(Section~\ref{sec:types:variable_declaration}). A parameter that must exist
without being read, because every call passes an argument for it, is named
\token{\_}, and the value passed for it is ignored. A name that starts and ends
with an underscore, such as \token{\_width\_}, is not checked either, but it
still says what the argument is, and the body can still read it. Both forms
work for any variable (Section~\ref{sec:memory:unused_variables}).

\subsection{Matching the arguments}

A call must give exactly one argument per parameter, and each argument must have
the type of its parameter. The call converts nothing. A literal such as
\token{12} takes the type the parameter asks for, but a variable must already
have that type.

\begin{lstlisting}[style=coloredverbatimError, escapechar=@, caption=Arguments that do not match the parameters]
use std::io;

fn line(width: i32) {
  for _ in 0 .. width {
    print("*");
  }
  println();
}

fn main() {
  line(12);    // ok, the literal is an i32
  let n = 12u8;
  line(@\hb{n}@);     // error, n is a u8
  line(@\hb{3, 4}@);  // error, one argument too many
}
\end{lstlisting}

Both calls are refused with the same message, \errmsg{E4226}{the call operator
  is not defined for \errhole{} and \errhole{}}, the holes being the function
and the arguments of the call. Under it, the compiler lists each function that
the call could have meant, and why it does not fit:
\errmsg{E4095}{incompatible types i32 and u8} for the first,
\errmsg{E4212}{1 parameters is expected, but 2 are provided} for the second.

\vfill\pagebreak
\section{Returning a value}
\label{sec:functions:return}

The functions of the previous section do something: they print. Many functions
compute something instead, and give the result back to their caller, which
decides what to do with it. Such a function declares the type of its result
after its parameters, with \token{->} followed by the type. The result is the
value of its body: Section~\ref{sec:control:blocks} showed that a block whose
last expression has no semicolon takes the value of that expression, and a
function body is a block like any other.

\lstinputlisting[style=coloredverbatim, caption={A function that returns a value, \textit{examples/functions/squares.yr}}, label=lst:functions:squares]{examples/functions/squares.yr}
\vspace{-10pt}%
\begin{lstlisting}[style=bashVerb, escapechar=@+]
@+\textcolor{teal!80}{alice@dev:\mtilde}@+$ ./squares
a = 3
b = 4
a * a + b * b = 25
and its square is 625
\end{lstlisting}

\token{square} takes an \token{i32} and returns an \token{i32}, the value of
\token{x * x}. A call to a function that returns a value is an expression, which
has that value, so it can go wherever a value of its type can go: into a variable
on line~10, into an addition on the same line, or as the argument of a call on
line~12, where the value of \token{square(sum)} is printed.

A function that declares no result, like \token{separator}, returns nothing. Its
result type is \token{void}, and \token{fn separator()} could be written
\token{fn separator()-> void}. Its call is not a value, and cannot be stored in a
variable.

\subsection{The trailing semicolon}

The rule of Section~\ref{sec:control:blocks} applies here too: a semicolon after
the last expression turns it into a statement, and the block then has no value.
The body of the next \token{square} computes \token{x * x}, throws it away, and
gives no value, while the declaration promises an \token{i32}.

\begin{lstlisting}[style=coloredverbatimError, escapechar=@, caption=A body with no value]
use std::io;

fn square(x: i32)-> @\hb{i32}@ {
  x * x@\hb{;}@ // error
}

fn main() {
  println(square(3));
}
\end{lstlisting}

The compiler reports \errmsg{E4095}{incompatible types i32 and void} and underlines the
promised \token{i32}, although the mistake is the semicolon two lines below: the
body's type, \token{void}, does not match the declared one.

\subsection{Leaving early: \texttt{return}}

The value of the body is the result of a function that runs to its end. The
keyword \token{return} leaves the function at once, from anywhere in the body,
with the value written after it. A \token{return} inside a loop leaves the loop
and the function together.

\lstinputlisting[style=coloredverbatim, caption={A function that returns early, \textit{examples/functions/primes.yr}}, label=lst:functions:primes]{examples/functions/primes.yr}
\vspace{-10pt}%
\begin{lstlisting}[style=bashVerb]
2 3 5 7 11 13 17 19 23 29 31 37 41 43 47
\end{lstlisting}

\token{isPrime} answers a question, so it returns a \token{bool}, and
\token{main} uses the answer directly as the condition of an \token{if}. The
function reads in the order in which the question is settled:

\begin{itemize}
\item a number below 2 is not prime, and line~5 says so without going further;
\item otherwise, the loop tries each divisor from 2 to \token{n - 1}, and the
  first one that divides \token{n} settles the question on line~9: the loop and
  the function stop there;
\item a number that reaches line~12 has no divisor, and the value of the body,
  \token{true}, is the answer.
\end{itemize}

Written without \token{return}, the same test needs a mutable variable to hold
the answer and a \token{break} to leave the loop, as in the exercises of
Chapter~\ref{chap:control}. With it, each case is answered where it is found.

A \token{return} in a function that returns nothing is written without a value:
\token{return;}. Since \token{return} leaves the function, a statement written
right after it could never run. The compiler refuses it, with the message
\errmsg{E4260}{unreachable statement}.

\begin{lstlisting}[style=coloredverbatimError, escapechar=@, caption=A statement after \token{return}]
use std::io;

fn half(n: i32)-> i32 {
  return n / 2;
  @\hb{println("Halved.");}@ // error
}

fn main() {
  println(half(10));
}
\end{lstlisting}

\vfill\pagebreak
\section{Naming the arguments}
\label{sec:functions:named_arguments}

In a call such as \token{printDate(4, 7, 1776)}, nothing says which number is
the day and which the month. The answer depends on the declaration of
\token{printDate}, and a reader used to the other convention reads the fourth of
July as the seventh of April. Two arguments of the same type can be swapped by
mistake, and the compiler cannot notice, since both orders have the right types.

\subsection{Named arguments}

An argument can name the parameter it is for, with the name of the parameter,
the arrow \token{->}, and the value: \token{printDate(day-> 4, month-> 7,
year-> 1776)}. A named argument goes to the parameter it names, wherever it is
written in the call, so the named arguments can come in any order. The arguments
without a name, called \textit{positional}, fill the remaining parameters in
order. Line~10 of the next listing names the year and the month, and the
positional \token{4} goes to the one parameter left, \token{day}.

\lstinputlisting[style=coloredverbatim, caption={Named arguments and a default value, \textit{examples/functions/dates.yr}}, label=lst:functions:dates]{examples/functions/dates.yr}
\vspace{-10pt}%
\begin{lstlisting}[style=bashVerb]
4/7/1776
4/7/1776
4/7/1776
14-7-1789
14.7.1789
\end{lstlisting}

The first three calls give the same arguments to the same parameters, and print
the same line. A parameter must get exactly one argument, so naming one twice,
or naming one that a positional argument has already filled, is refused.

\subsection{Default values}

A parameter can have a \textit{default value}, written after its type with
\token{=}, like the initial value of a variable. A call that gives no argument
for that parameter gets the default. \token{separator}, on line~3, defaults to
\token{'/'}: the calls of lines~8 to~10 do not mention it, and their dates are
printed with slashes.

A call that wants another value gives it \emph{by name}, as on lines~11 and~12.
A parameter with a default value is never filled by a positional argument, so the
last parameter of a declaration can have a default without making the others
optional, and a date is never read as a separator by accident.

\begin{lstlisting}[style=coloredverbatimError, escapechar=@, caption=A default value given by position]
use std::io;

fn printDate(day: i32, month: i32, year: i32, separator: c8 = '/') {
  println(day, separator, month, separator, year);
}

fn main() {
  printDate(14, 7, 1789, @\hb{'-'}@); // error
}
\end{lstlisting}

The compiler lists two reasons: \errmsg{E4212}{3 parameters are expected, but 4
  are provided}, since only the parameters without a default are filled by
position, and \errmsg{E4283}{the parameter separator has a default value, it can
  only be passed by name}.

\subsection{Arguments that must be named}

A named argument protects a call only if its author thinks of naming it. A
function can insist, by writing \token{@} in front of a parameter's name: that
parameter can then only be given by name. Declared with
\token{@day}, \token{@month} and \token{@year}, \token{printDate} can no longer
be called with three bare numbers, and every call says which date it means.

\begin{lstlisting}[style=coloredverbatimError, escapechar=|, caption=Parameters that must be named]
use std::io;

fn printDate(@day: i32, @month: i32, @year: i32) {
  println(day, "/", month, "/", year);
}

fn main() {
  printDate(day-> 14, month-> 7, year-> 1789); // ok
  printDate(|\hb{14, 7, 1789}|);                       // error
}
\end{lstlisting}

The second call is refused. The compiler counts
\errmsg{E4212}{0 parameters are expected, but 3 are provided}, since no parameter
can be filled by position, then gives one line per parameter:
\errmsg{E4284}{the parameter day can only be passed by name}, and the same for
\token{month} and \token{year}. The
\token{@} belongs to the declaration only; the body names the parameter
\token{day}, without it.

Whether a parameter should be named is a question about its calls. When the
function has a single parameter, or when the order of the arguments is obvious,
as in \token{rectangle(width, height)}, names only make the calls longer. When
several parameters have the same type and no natural order, as a day, a month
and a year, \token{@} rules out a whole class of mistakes. The calendar of
Section~\ref{sec:functions:calendar} uses it for exactly that.

\subsection{Reading, explained}
\label{sec:functions:read}

Chapter~\ref{chap:control} asked for \token{read!i32(ask-> "Your guess: ")} to
be taken as a fixed formula. Its second half is now ordinary. \token{read} is a
function of \token{std::io}, and its declaration is, in essence:

%% check: skip -- the real declaration needs templates; this is its shape
\begin{lstlisting}[style=coloredverbatim]
fn read(ask: [c8] = "")-> T
\end{lstlisting}

\token{ask} is a parameter with a default value, the empty string, so
\token{read!i32()} prints nothing before waiting, and a call that wants a prompt
gives it by name, \token{ask-> "Your guess: "}. The \token{T} is the part that
stays unexplained: \token{read} can return a value of many types, and
\token{!i32} is how the call chooses one. Functions that work on any type are
called \textit{templates}, and have a chapter of their own,
Chapter~\ref{chap:templates}.

\subsection{Calling with a dot}
\label{sec:functions:ucs}

A call can also be written with its first argument in front of the function
name, followed by a dot: \token{n.square()} is the call \token{square(n)}, and
\token{a.add(b)} is \token{add(a, b)}. The compiler rewrites the first form into
the second, so it works with any function, without changing its declaration.
This way of calling is the \textit{uniform call syntax}.

\begin{lstlisting}[style=coloredverbatim, caption=Calls written with a dot, label=lst:functions:ucs]
use std::io;

fn square(x: i32)-> i32 {
  x * x
}

fn add(a: i32, b: i32)-> i32 {
  a + b
}

fn main() {
  let n = 3;
  println(square(add(square(n), 1)));
  println(n.square().add(1).square());
}
\end{lstlisting}
\vspace{-10pt}%
\begin{lstlisting}[style=bashVerb]
100
100
\end{lstlisting}

Both lines compute the same thing: square 3, add 1, and square the result.
Nested calls, on line~13, are read from the inside out, starting with the
innermost \token{square(n)}. The chain on line~14 is read from left to right, in
the order the operations happen. The dot works after a literal as well, with one
catch: \token{3.square()} is read as the floating point number \token{3.} followed
by a name, and the literal must be put between parentheses,
\token{(3).square()}.

\vfill\pagebreak
\section{Functions that call themselves}
\label{sec:functions:recursion}

A function can call any function, including itself. Such a function is
\textit{recursive}. It solves a problem by solving a smaller version of the same
problem, and it relies on the smallest version having an answer that needs no
further call. A countdown from 3 is ``say 3, then count down from 2'', and a
countdown from 0 is simply ``liftoff''.

\lstinputlisting[style=coloredverbatim, caption={A recursive countdown, \textit{examples/functions/countdown.yr}}, label=lst:functions:countdown]{examples/functions/countdown.yr}
\vspace{-10pt}%
\begin{lstlisting}[style=bashVerb]
3...
2...
1...
Liftoff!
\end{lstlisting}

Every recursive function has the same two parts, and \token{countdown} shows
both:

\begin{itemize}
\item a \textit{base case}, the problem small enough to answer directly, without
  calling the function again: here \token{n == 0}, on line~4;
\item a \textit{recursive case}, which does a little of the work and hands the
  rest to a call on a smaller problem: here, print \token{n}, then count down from
  \token{n - 1}, on line~8.
\end{itemize}

Each call has its own parameter. When \token{countdown(3)} calls
\token{countdown(2)}, the new call gets a new \token{n}, holding 2, and the
\token{n} of the first call still holds 3. Nothing is shared between them.

\subsection{Waiting for the answer}

In \token{countdown}, the recursive call is the last thing the function does. In
most recursive functions, the call's result is needed to finish the work. The
\textit{factorial} of a number $n$, written $n!$, is the product
$1 \times 2 \times \dots \times n$, with $0! = 1$. It is also $n$ times the
factorial of $n - 1$, and that sentence is a recursive function.

\lstinputlisting[style=coloredverbatim, caption={The factorial, \textit{examples/functions/factorial.yr}}, label=lst:functions:factorial]{examples/functions/factorial.yr}
\vspace{-10pt}%
\begin{lstlisting}[style=bashVerb, escapechar=@+]
@+\textcolor{teal!80}{alice@dev:\mtilde}@+$ ./factorial
n = 20
20! = 2432902008176640000
\end{lstlisting}

To compute \token{n * factorial(n - 1)} on line~7, a call must wait for the
result of the call it makes, which waits for the next one, and so on down to the
base case. Figure~\ref{fig:functions:factorial_calls} shows the calls made by
\token{factorial(3)}. The chain goes down until \token{n} is 0, whose call
answers 1 at once. Then the results travel back up, and each waiting call
multiplies the result it receives by its own \token{n}.

\begin{figure}[h]
  \centering
  \begin{adjustbox}{max width=\linewidth}
    \begin{tikzpicture}[node distance=10mm,
        call/.style={cfstep, minimum width=52mm, align=left, text width=48mm}]

      \node[cfterm] (main) {main: println(factorial(3))};
      \node[call, below=of main] (f3) {n = 3:\quad 3 * factorial(2)};
      \node[call, below=of f3] (f2) {n = 2:\quad 2 * factorial(1)};
      \node[call, below=of f2] (f1) {n = 1:\quad 1 * factorial(0)};
      \node[call, below=of f1] (f0) {n = 0:\quad 1};

      \foreach \upper/\lower in {main/f3, f3/f2, f2/f1, f1/f0} {
        \draw[cfflow] ($(\lower.north west) + (8mm, 10mm)$) --
          node[cflab, left] {calls} ($(\lower.north west) + (8mm, 0)$);
      }
      \foreach \upper/\lower/\value in {main/f3/6, f3/f2/2, f2/f1/1, f1/f0/1} {
        \draw[cfflow, dashed] ($(\lower.north east) + (-8mm, 0)$) --
          node[cflab, right] {returns \value} ($(\lower.north east) + (-8mm, 10mm)$);
      }
    \end{tikzpicture}
  \end{adjustbox}
  \caption{\label{fig:functions:factorial_calls} The calls made by
    \token{factorial(3)}. Each call waits for the one below it, and each has its
    own \token{n}. The call with \token{n = 0} answers without calling, and the
    results travel back up, each call multiplying by its own \token{n}.}
\end{figure}


The factorial grows fast. \token{factorial(20)} is the largest that fits in an
\token{i64}; \token{factorial(21)} overflows, and the program prints a negative
number, for the reasons of Section~\ref{sec:types:integer_overflow}.

Two functions can also call each other. The next two decide whether a positive
number is even or odd, knowing only that 0 is even: a number is even when the
number before it is odd, and odd when the number before it is even. Each
function is the recursive case of the other.

\begin{lstlisting}[style=coloredverbatim, caption=Two functions calling each other, label=lst:functions:parity]
use std::io;

fn isEven(n: i32)-> bool {
  if n == 0 { true } else { isOdd(n - 1) }
}

fn isOdd(n: i32)-> bool {
  if n == 0 { false } else { isEven(n - 1) }
}

fn main() {
  println(isEven(10), " ", isOdd(10));
}
\end{lstlisting}
\vspace{-10pt}%
\begin{lstlisting}[style=bashVerb]
true false
\end{lstlisting}

This is a slow way to test parity, which \token{n \% 2 == 0} answers in one
operation. It shows that the base case can live in any of the functions of the
cycle, as long as every chain of calls reaches one.

\subsection{Infinite recursion}
\label{sec:functions:infinite_recursion}

A loop whose condition never becomes false runs forever. A recursion whose base
case is never reached does worse: it crashes. The next countdown has lost its
base case, and every call makes another.

\lstinputlisting[style=coloredverbatim, caption={A recursion with no base case, \textit{examples/functions/runaway.yr}}, label=lst:functions:runaway]{examples/functions/runaway.yr}
\vspace{-10pt}%
\begin{lstlisting}[style=bashVerb, escapechar=@+]
@+\textcolor{teal!80}{alice@dev:\mtilde}@+$ ./runaway
3...
2...
1...
0...
-1...
@+\textit{(about 130\,000 more lines)}@+
-130822...
-130823...
Segmentation fault (core dumped)
\end{lstlisting}

The compiler accepts this program, and it seems at first to count forever, into
the negative numbers. Then it stops, and the last line, printed by the terminal
rather than by the program, says that the operating system killed it.

The reason is the waiting that Figure~\ref{fig:functions:factorial_calls} shows.
A call that has not returned yet keeps its parameters and its place in the
caller somewhere in memory, in a region called the \textit{call stack}. The
stack has a fixed size, a few megabytes, and each call takes a little of it.
Here no call ever returns, so the calls pile up until the stack is full, and the
next call has nowhere to go. This is a \textit{stack overflow}. The number of
calls it takes depends on the machine; about 130\,000 was the count on the
machine that ran this listing. Chapter~\ref{chap:layout} explains what the stack
holds.

A base case that exists but is never reached has the same effect. With the base
case of Listing~\ref{lst:functions:countdown}, the call \token{countdown(-1)}
tests \token{-1 == 0}, then \token{-2 == 0}, and so on: \token{n} moves away from
0 and never meets it. Writing the test as \token{n <= 0} closes the gap, and a
recursive function is safest when its base case catches every value that it
cannot reduce.

Two questions check a recursive function against both mistakes:

\begin{itemize}
\item Is there a case that returns without calling the function?
\item Does every recursive call get a problem strictly closer to that case,
  whatever the arguments?
\end{itemize}

\subsection{Too deep}

A recursion can also overflow the stack while being correct. The sum of the
numbers from 1 to $n$ is $n$ plus the sum from 1 to $n - 1$, and a function
written that way gives the right answer, as long as it is not asked for too
much.

\lstinputlisting[style=coloredverbatim, caption={A correct recursion that goes too deep, \textit{examples/functions/deep\_sum.yr}}, label=lst:functions:deep_sum]{examples/functions/deep_sum.yr}
\vspace{-10pt}%
\begin{lstlisting}[style=bashVerb, escapechar=@+]
@+\textcolor{teal!80}{alice@dev:\mtilde}@+$ ./deep_sum
500500
Segmentation fault (core dumped)
\end{lstlisting}

The sum up to 1\,000 needs 1\,001 calls waiting at the same time, and fits. The
sum up to 100\,000 needs 100\,001 of them, and overflows the stack before the
base case is reached. The loop of Listing~\ref{lst:control:for_sum} computes the
same sum for any $n$, with one call and one variable.

A recursion is a good fit when each call does something the loop version would
have to organize by hand, as in the next subsection or in the exercises. When
the recursion only counts down, one step at a time, a loop does the same work
without piling up calls.

\subsection{The Ackermann function}
\label{sec:functions:ackermann}

The Ackermann function, named after the mathematician Wilhelm Ackermann, takes
two numbers $m$ and $n$ and is defined by three rules:

\begin{itemize}
\item when $m$ is 0, the answer is $n + 1$;
\item when $n$ is 0, the answer is the value for $m - 1$ and $1$;
\item otherwise, the answer is the value for $m - 1$ and the value for $m$ and
  $n - 1$.
\end{itemize}

Each rule becomes one branch of the function. The third is unusual: the second
argument of the recursive call is itself a recursive call.

\lstinputlisting[style=coloredverbatim, caption={The Ackermann function, \textit{examples/functions/ackermann.yr}}, label=lst:functions:ackermann]{examples/functions/ackermann.yr}
\vspace{-10pt}%
\begin{lstlisting}[style=bashVerb]
1 2 3 4 5 6
2 3 4 5 6 7
3 5 7 9 11 13
5 13 29 61 125 253
\end{lstlisting}

The function always terminates. Every call either lowers $m$, or keeps $m$ and
lowers $n$, so the calls always reach $m = 0$ in the end. Yet the table shows
how fast the values grow. The row $m = 1$ is $n + 2$, the row $m = 2$ is
$2n + 3$, and the row $m = 3$ is already $2^{n+3} - 3$, which doubles at each
step. The row $m = 4$ is worse still:

\begin{itemize}
\item \token{ackermann(4, 1)} is 65\,533. Computing it takes almost three billion
  calls, some ten seconds on the machine that ran these listings, with 65\,536
  calls waiting at the deepest point.
\item \token{ackermann(4, 2)} is $2^{65536} - 3$, a number with 19\,729 digits.
  It does not fit in any integer type, and no computer could make the number of
  calls this function needs to compute it.
\end{itemize}

The Ackermann function is famous for a reason beyond its growth. Mathematicians
proved that it cannot be computed with \token{for} loops alone, whose number of
iterations is known when they start: recursion, or a loop whose end is not known
in advance, is needed. It is also a reminder that a function can be correct,
terminate, and still never finish in practice. How many calls a function makes
matters as much as whether it stops.

\vfill\pagebreak
\section{Several functions with the same name}
\label{sec:functions:overloading}

The area of a square needs one length, and the area of a rectangle two. Both
functions compute an area, and names such as \token{squareArea} and
\token{rectangleArea} only repeat what the arguments already say. Ymir lets
several functions share a name, as long as their parameters differ, in number or
in type. The name is then \textit{overloaded}, and each call is sent to the
function whose parameters match its arguments.

\lstinputlisting[style=coloredverbatim, caption={Two functions named \token{area}, \textit{examples/functions/area.yr}}, label=lst:functions:area]{examples/functions/area.yr}
\vspace{-10pt}%
\begin{lstlisting}[style=bashVerb]
A 3 by 3 square: 9
A 2 by 5 rectangle: 10
\end{lstlisting}

The call on line~12 has one argument, and only the first \token{area} takes one.
The call on line~13 has two, and goes to the second. The choice is made by the
compiler, once, when it reads the call, not each time the program runs.

Functions with the same number of parameters are told apart by their types. The
three functions of the next listing take one argument each, of three different
types, and each call goes to the function whose parameter has the type of its
argument.

\lstinputlisting[style=coloredverbatim, caption={Overloading on the type of the argument, \textit{examples/functions/describe.yr}}, label=lst:functions:describe]{examples/functions/describe.yr}
\vspace{-10pt}%
\begin{lstlisting}[style=bashVerb]
42 is an integer
2.5 is a floating point number
true is a boolean
\end{lstlisting}

\token{println} works the same way from the caller's point of view: it accepts
integers, floats, strings and Booleans alike, and prints each the way its type
requires. How it is written is the subject of Chapter~\ref{chap:templates}.

\subsection{What must differ}

The compiler chooses a function from the arguments of the call, and from nothing
else. Two functions it could not tell apart from their arguments are refused.
The names of the parameters, which the call does not have to write, do not tell
two functions apart.

\begin{lstlisting}[style=coloredverbatimError, escapechar=@, caption=Two functions that differ only by the names of their parameters]
use std::io;

fn area(side: f64)-> f64 {
  side * side
}

fn @\hb{area}@(width: f64)-> f64 { // error
  width * width
}

fn main() {
  println(area(3.0));
}
\end{lstlisting}

The compiler reports \errmsg{E4023}{colliding function definitions \errhole{} and
  \errhole{}}, naming both. The result
type does not tell two functions apart either, since a call does not say which
type it expects.

\begin{lstlisting}[style=coloredverbatimError, escapechar=@, caption=Two functions that differ only by their result]
use std::io;

fn zero()-> i32 {
  0
}

fn @\hb{zero}@()-> f64 { // error
  0.0
}

fn main() {
  println(@\hb{zero()}@); // error
}
\end{lstlisting}

The two declarations collide, as before. The call is reported too: nothing in
\token{zero()} says whether it wants the integer or the floating point zero, and
the compiler says \errmsg{E4280}{\errhole{} called with \{\} works with multiple
  candidates}, the hole being \token{zero}.

\tipbox{Overloading is for functions that do the same job on different inputs,
  so that the reader of a call knows what it does without knowing which of them
  runs. Two functions that do different things deserve different names.}

\vfill\pagebreak
\section{Putting it together: a calendar}
\label{sec:functions:calendar}

The program of this section prints the calendar of a month, or of a whole year,
laid out week by week like the calendar on a wall. It is the first program of the
book long enough to need functions, and it is a useful one: the calendar of any
month from 1900 on, in a terminal.

\subsection{Splitting the problem}

A program of this size is not written from the first line to the last. It is
easier to start from what it must do, and to split that into questions small
enough to be answered by one short function each. Printing the calendar of a
month raises these questions:

\begin{itemize}
\item \textit{What is the month called?} A \token{match} on the month number
  answers it.
\item \textit{How many days does the month have?} Exercise~6 of
  Chapter~\ref{chap:control} answered it for ordinary years. February has 29 days
  in a leap year, which raises the next question.
\item \textit{Is the year a leap year?} It is when it is divisible by 4, except
  for the years divisible by 100, which are leap years only when they are also
  divisible by 400: 2024 and 2000 are leap years, 1900 is not.
\item \textit{On which day of the week does the month start?} This decides how
  far to the right the 1st is printed. The program knows that January 1st, 1900
  was a Monday. It counts the days between that date and the 1st of the month,
  using the number of days in each year and in each month on the way, and the
  remainder of their division by 7 is the day of the week.
\end{itemize}

Each question becomes a function, and the functions that answer the last
questions call those that answer the first ones. Printing the calendar then
takes only a few lines, because the hard parts have names.

\subsection{The program}

\lstinputlisting[style=coloredverbatim, caption={A calendar, \textit{examples/functions/calendar.yr}}, label=lst:functions:calendar]{examples/functions/calendar.yr}
\vspace{-10pt}%
\begin{lstlisting}[style=bashVerb, escapechar=@+]
@+\textcolor{teal!80}{alice@dev:\mtilde}@+$ gyc calendar.yr -o calendar
@+\textcolor{teal!80}{alice@dev:\mtilde}@+$ ./calendar
Year: 1850
Please type a number from 1900 to 9999.
Year: 2026
Month (0 for the whole year): 9

     September 2026
Mo Tu We Th Fr Sa Su
    1  2  3  4  5  6
 7  8  9 10 11 12 13
14 15 16 17 18 19 20
21 22 23 24 25 26 27
28 29 30
\end{lstlisting}

\subsection{How it works}

The functions are declared in the order of the questions, each after the ones it
calls. Nothing requires that order, but it lets the file be read from top to
bottom.

\begin{itemize}
\item \token{isLeapYear}, on line~3, is one Boolean expression, whose value is
  the answer.
\item \token{daysIn} is overloaded. With one argument, on line~7, it gives the
  number of days in a year; with two, on line~11, the number of days in a month
  of a given year. Both are ``the number of days in'' something, and the
  arguments say what. The arm for February is an \token{if} used as a value.
\item \token{weekday}, on line~38, counts the days from January 1st, 1900. The
  first loop adds the days of every year before \token{year}, the second the days
  of every month before \token{month}, and line~46 adds the days of the month
  before \token{day}. Its three parameters have the same type and a day, a month
  and a year are written in different orders in different countries, so all
  three are marked \token{@}. The call on line~59 cannot mix them up.
\item \token{printDay}, on line~49, prints a day on two characters, padded with a
  space, so that the columns line up.
\item \token{printCalendar}, on line~56, prints the title and the header, then
  three spaces for each day of the week before the 1st, then the days. After a
  day, the number of cells printed on the current row is \token{first + day},
  counting the blank ones; when it is a multiple of 7, the week is complete and
  the row ends. The \token{if} on line~70 ends the last row, unless the month
  ended on a Sunday and the row is already ended.
\item \token{printCalendar} is overloaded too. The version with one parameter,
  on line~75, prints a whole year by calling the other for each month.
\item \token{askNumber}, on line~82, asks the same question until the answer is
  within bounds. Its bounds must be named, so the call on line~93 reads as
  ``a number from 1900 to 9999''. The \token{loop} produces the number accepted,
  with \token{break n} (Section~\ref{sec:control:loop}).
\item \token{main} is left with nothing to do but ask and choose.
\end{itemize}

Every function of the program is short enough to be checked on its own. If a
month started on the wrong day, only \token{weekday} could be at fault, and it
can be tested alone, with a few dates whose day is known. The last exercise of
the chapter reuses it.

\vfill%
\pagebreak

\section*{Exercises}
\markright{\MakeUppercase{Exercises}}%
\subsection*{Exercise 1}

Write a function \token{smallest} that returns the smaller of two \token{i32}.
Then overload it with a version that takes three, written with calls to the
first, and no \token{if} of its own.

\subsection*{Exercise 2}

What does the following program print? Work it out by hand, then check your
answer by compiling and running it.

\begin{lstlisting}[style=coloredverbatim]
use std::io;

fn show(label: [c8], value: i32, unit: [c8] = "") {
  println(label, ": ", value, unit);
}

fn twice(x: i32)-> i32 {
  x * 2
}

fn main() {
  let n = 5;
  show("n", n);
  show("twice", n.twice(), unit-> " points");
  show(value-> n.twice().twice(), "four times");
  show(unit-> "!", "sum", twice(n) + n);
}
\end{lstlisting}

\subsection*{Exercise 3}

The following program should print the average of 4 and 10, kept within the
bounds 0 and 5. It contains three errors. Find them, correct them, and compile the
program to test your corrections.

\begin{lstlisting}[style=coloredverbatimError, caption=Ymir program with errors]
use std::io;

fn average(a: i32, b: i32)-> i32 {
  (a + b) / 2;
}

fn clamp(value: i32, @low: i32, @high: i32)-> i32 {
  if value < low {
    value = low;
  }
  if value > high { high } else { value }
}

fn main() {
  let m = average(4, 10);
  println(clamp(m, 0, 5));
}
\end{lstlisting}

\subsection*{Exercise 4}

Write a recursive function \token{power(base: i64, exponent: i32)-> i64} that
computes \token{base} to the power \token{exponent}, knowing that any number to
the power 0 is 1, and that $b^e = b \times b^{e - 1}$. What happens if it is
called with a negative exponent? Change the function so that it returns 1
instead.

\subsection*{Exercise 5}

Exercise~4 of Chapter~\ref{chap:control} computed the sum of the digits of a
number with a loop. Write it as a recursive function instead. (\textit{Hint: a
number below 10 is its own sum of digits; for a larger one, what is the relation
between the sum of the digits of \token{n} and that of \token{n / 10}?})

\subsection*{Exercise 6}

The Fibonacci sequence starts with 0 and 1, and each number after them is the sum
of the two before it: 0, 1, 1, 2, 3, 5, 8, 13, \dots\ Write a recursive function
\token{fibonacci(n: i64)-> i64} that returns the number at position \token{n},
counting from 0, straight from that definition. Run it for 30, 40 and 45. Why
does it get so slow, when a loop adds up millions of numbers in the blink of an
eye? Write a version with a loop that answers at once, even for 90.

\subsection*{Exercise 7}

The greatest common divisor of two numbers is the largest number that divides
both. Euclid found, some 2\,300 years ago, that the greatest common divisor of
\token{a} and 0 is \token{a}, and that the greatest common divisor of \token{a}
and \token{b} is that of \token{b} and \token{a \% b}. Write a recursive function
\token{gcd} from these two rules, and check that the greatest common divisor of
1071 and 462 is 21. Why does the recursion always end?

\subsection*{Exercise 8}

Following the rules of Section~\ref{sec:functions:ackermann}, compute
\token{ackermann(1, 2)} by hand, writing each call that the function makes. How
many calls are there in all?

\subsection*{Exercise 9}

Write a program that asks for a date and prints the day of the week it fell on,
such as \textit{That day is a Sunday.} for July 20th, 1969. Reuse the functions of
Listing~\ref{lst:functions:calendar} that you need, and make sure that the day
typed exists in the month typed.

\vfill%
\pagebreak

\forceevenpage{}

\section*{Solutions}
\markright{\MakeUppercase{Solutions}}%
\subsection*{Exercise 1}

The smallest of three numbers is the smallest of the third and of the smaller of
the first two. The version with three parameters calls the one with two, twice.

\lstinputlisting[style=coloredverbatim, caption={Solution for exercise 1, \textit{examples/functions/solutions/smallest.yr}}]{examples/functions/solutions/smallest.yr}

\subsection*{Exercise 2}

\begin{lstlisting}[style=bashVerb]
n: 5
twice: 10 points
four times: 20
sum: 15!
\end{lstlisting}

The first call leaves \token{unit} at its default, the empty string. The third
names \token{value}, so the positional \token{"four times"} goes to the first
free parameter, \token{label}. In the last, \token{unit} is named first, and the
two positional arguments fill \token{label} and \token{value} in order.

\subsection*{Exercise 3}

The body of \token{average} ends with a semicolon, so it has no value.
\token{value} is a parameter, and cannot be assigned: the three cases are
written as one \token{if} used as a value. And \token{low} and \token{high} are
marked \token{@}, so the call must name them. The program prints 5: the average
is 7, above the upper bound.

\begin{lstlisting}[style=coloredverbatim, escapechar=|, caption=Solution for exercise 3]
use std::io;

fn average(a: i32, b: i32)-> i32 {
  |\hcb{(a + b) / 2}|
}

fn clamp(value: i32, @low: i32, @high: i32)-> i32 {
  if value < low {
    |\hcb{low}|
  } else if value > high {
    high
  } else {
    value
  }
}

fn main() {
  let m = average(4, 10);
  println(clamp(m, |\hcb{low-> 0, high-> 5}|));
}
\end{lstlisting}

\subsection*{Exercise 4}

With the base case \token{exponent == 0}, a negative exponent moves away from 0
at each call, and the recursion never ends: the program crashes with a stack
overflow, as in Section~\ref{sec:functions:infinite_recursion}. Testing
\token{exponent <= 0} catches it.

\lstinputlisting[style=coloredverbatim, caption={Solution for exercise 4, \textit{examples/functions/solutions/power.yr}}]{examples/functions/solutions/power.yr}
\vspace{-10pt}%
\begin{lstlisting}[style=bashVerb]
1024
81
1
1
\end{lstlisting}

\subsection*{Exercise 5}

The sum of the digits of \token{n} is its last digit, \token{n \% 10}, plus the
sum of the digits of the others, \token{n / 10}.

\lstinputlisting[style=coloredverbatim, caption={Solution for exercise 5, \textit{examples/functions/solutions/digits.yr}}]{examples/functions/solutions/digits.yr}

\subsection*{Exercise 6}

\lstinputlisting[style=coloredverbatim, caption={Solution for exercise 6, recursive, \textit{examples/functions/solutions/fibonacci.yr}}]{examples/functions/solutions/fibonacci.yr}

\token{fibonacci(30)} answers at once, \token{fibonacci(40)} takes most of a
second, and \token{fibonacci(45)} several seconds. Each call makes two more, and
they compute the same values again and again: \token{fibonacci(40)} calls
\token{fibonacci(39)} and \token{fibonacci(38)}, but \token{fibonacci(39)}
computes \token{fibonacci(38)} a second time, and each of them recomputes all the
smaller ones. In all, \token{fibonacci(40)} makes more than 331 million calls,
and each step of 5 in \token{n} multiplies that by about 11.

The loop keeps the last two numbers of the sequence and moves them forward
\token{n} times, so it makes \token{n} additions in all.

\lstinputlisting[style=coloredverbatim, caption={Solution for exercise 6, with a loop, \textit{examples/functions/solutions/fibonacci\_loop.yr}}]{examples/functions/solutions/fibonacci_loop.yr}
\vspace{-10pt}%
\begin{lstlisting}[style=bashVerb, escapechar=@+]
@+\textcolor{teal!80}{alice@dev:\mtilde}@+$ ./fibonacci_loop
n = 90
fibonacci(90) = 2880067194370816120
\end{lstlisting}

\subsection*{Exercise 7}

\lstinputlisting[style=coloredverbatim, caption={Solution for exercise 7, \textit{examples/functions/solutions/gcd.yr}}]{examples/functions/solutions/gcd.yr}
\vspace{-10pt}%
\begin{lstlisting}[style=bashVerb, escapechar=@+]
@+\textcolor{teal!80}{alice@dev:\mtilde}@+$ ./gcd
a = 1071
b = 462
gcd(1071, 462) = 21
\end{lstlisting}

The remainder of a division by \token{b} is smaller than \token{b}, so for
positive numbers the second argument gets smaller at each call, and eventually
reaches the base case, 0.

\subsection*{Exercise 8}

Each line is one call. The calls with $m = 1$ wait for the value of their inner
call, then make one more call with $m = 0$.

\begin{lstlisting}[style=bashVerb]
ackermann(1, 2) = ackermann(0, ackermann(1, 1))
ackermann(1, 1) = ackermann(0, ackermann(1, 0))
ackermann(1, 0) = ackermann(0, 1)
ackermann(0, 1) = 2
ackermann(0, 2) = 3          so ackermann(1, 1) = 3
ackermann(0, 3) = 4          so ackermann(1, 2) = 4
\end{lstlisting}

There are 6 calls in all, three with $m = 1$ and three with $m = 0$. The value,
4, is the one in the table of Listing~\ref{lst:functions:ackermann}.

\subsection*{Exercise 9}

\token{weekday} and the functions it calls are copied unchanged. The upper bound
of the day depends on the month and the year, so it is computed by
\token{daysIn}, and the numbers from 0 to 6 are turned into names by a
\token{match}, as \token{monthName} does for months.

\lstinputlisting[style=coloredverbatim, caption={Solution for exercise 9, \textit{examples/functions/solutions/weekday.yr}}]{examples/functions/solutions/weekday.yr}
\vspace{-10pt}%
\begin{lstlisting}[style=bashVerb, escapechar=@+]
@+\textcolor{teal!80}{alice@dev:\mtilde}@+$ ./weekday
Year: 1969
Month: 7
Day: 20
That day is a Sunday.
\end{lstlisting}

